\documentclass[11pt]{article}

\usepackage[T1]{fontenc}
\usepackage{lmodern}
\usepackage{amsmath,amssymb,amsthm,mathtools}
\usepackage{microtype}
\usepackage{etoolbox}
\usepackage[margin=1.08in]{geometry}
\usepackage{xcolor}
\usepackage[colorlinks=true,linkcolor=blue!55!black,citecolor=blue!55!black,
  urlcolor=blue!55!black]{hyperref}
\usepackage[nameinlink,noabbrev]{cleveref}

\newtheorem{theorem}{Theorem}[section]
\newtheorem{proposition}[theorem]{Proposition}
\newtheorem{corollary}[theorem]{Corollary}
\newtheorem{lemma}[theorem]{Lemma}
\theoremstyle{definition}
\newtheorem{definition}[theorem]{Definition}
\theoremstyle{remark}
\newtheorem{remark}[theorem]{Remark}

\BeforeBeginEnvironment{proof}{\smallskip}
\AfterEndEnvironment{proof}{\smallskip}

\newcommand{\C}{\mathbb C}
\newcommand{\N}{\mathbb N}
\newcommand{\PP}{\mathbb P}
\newcommand{\A}{\mathcal A}
\newcommand{\F}{\mathcal F}
\newcommand{\Efin}{\mathcal E_{\mathrm{fin}}}
\newcommand{\Pcal}{\mathcal P}
\newcommand{\K}{\mathcal K}
\newcommand{\J}{\mathcal J}
\newcommand{\ord}{\operatorname{ord}}
\newcommand{\Frac}{\operatorname{Frac}}
\newcommand{\supp}{\operatorname{supp}}
\newcommand{\trdeg}{\operatorname{trdeg}}
\newcommand{\dtrdeg}{\operatorname{dtrdeg}}
\newcommand{\Span}{\operatorname{span}}

\title{Prime-Profile Differential Fields and the\\
Valuation-Quotient Hypertranscendence Law}

\author{Michael Haufschild}
\date{}

\begin{document}
\maketitle

\begin{abstract}
Let $\A$ be the Dirichlet-convolution domain of complex arithmetic functions, let
$\F=\Frac(\A)$, and equip $\F$ with the logarithmic derivation
$(Da)(n)=\mathrm i\log(n)a(n)$.  A profile $u\in t\C[[t]]$ determines a prime-radial source
$a_u$: its value at $p^n$ is $[t^n]u$, independently of $p$, and it vanishes away from positive
prime powers.  We prove an exact source-to-field correspondence for arbitrary profile subspaces.
If $\Efin$ is the differential field of fraction-field elements with only finitely many active
prime coordinates and
\[
  \K(W)=\Efin\langle a_u:u\in W\rangle_D,
\]
then
\[
  a_u\in\K(W)\Longleftrightarrow u\in W,
  \qquad
  \K(U)\subseteq\K(W)\Longleftrightarrow U\subseteq W.
\]

The main result identifies every finite valuation quotient with an intrinsic block of relatively
hypertranscendental channels.  Let $\pi_N$ retain coefficients of degrees $1,\ldots,N$ and put
$W_{>N}=W\cap\ker\pi_N$.  If $b_1,\ldots,b_h$ lift a basis of $\pi_N(W)$, then they
differentially generate $\K(W)$ over the complete tail field $\K(W_{>N})$ and their carriers are
jointly differentially independent over that field.  Hence they form a relative differential
transcendence basis, and the ordinary field generated by their jets through order $r$ has the
intrinsic rank
\[
 \trdeg_{\K(W_{>N})}
 \K(W_{>N})\bigl(D^ka_{b_j}:1\le j\le h,\ 0\le k\le r\bigr)
   =(r+1)\dim_\C\pi_N(W).
\]
The field and the displayed degree are independent of the chosen quotient basis and its lifts.

Finally, raywise Cauchy fusion is identified with Dirichlet convolution followed by the canonical
projection to positive prime powers.  In the radial sector this becomes profile multiplication,
so fusion-closed radial fields correspond exactly to multiplicatively closed profile subspaces.
Order semigroups and conductor periodicity then appear as classical filtered shadows of the
tail-relative law, rather than as its source of novelty.
\end{abstract}

\medskip
\noindent\textbf{Keywords.}
hypertranscendence; differential algebra; arithmetic functions; Dirichlet convolution;
valuation filtration; formal power series.

\medskip
\noindent\textbf{2020 Mathematics Subject Classification.} Primary 12H05; Secondary 11J85,
11M41, 13A18.

\section{Introduction}

The derivative orbit of one element in a differential field is one-dimensional as a sequence but
potentially infinite-dimensional algebraically.  Prime-power arithmetic data carry a second
filtration: the valuation depth $n$ in $p^n$.  This paper determines how those two directions
interact for sources whose valuation profile is replicated on every prime ray.

The elementary encoding is as follows.  Put
\[
  \Pcal=t\C[[t]].
\]
For $u\in\Pcal$, define an arithmetic function $a_u$ by
\begin{equation}\label{eq:radial-source}
 a_u(p^n)=[t^n]u\quad(p\in\PP,\ n\ge1),
 \qquad
 a_u(m)=0\quad\text{if $m$ is not a positive prime power}.
\end{equation}
At the level of sources, this is merely a linear embedding.  The nontrivial issue is whether the
differential field generated by a family $(a_u)$ remembers the profile family or collapses
distinct valuation directions through nonlinear relations.

The answer is exact.  Inside the fraction field of the arithmetic-function ring we construct a
differential coefficient field $\Efin$ and prove that the assignment
\begin{equation}\label{eq:field-functor}
 W\longmapsto \K(W)=\Efin\langle a_u:u\in W\rangle_D
\end{equation}
is an order embedding from complex linear subspaces of $\Pcal$ to embedded differential fields.
This simultaneously classifies source membership, field containment, and equality in the radial
sector.

The main theorem is stronger than source recovery because it identifies the intrinsic extension
created across every valuation cutoff.  For $N\ge0$, let
\[
 \pi_N(u)=([t]u,\ldots,[t^N]u),
 \qquad W_{>N}=W\cap\ker\pi_N.
\]
The quotient $W/W_{>N}$ is canonically isomorphic to the finite-dimensional space $\pi_N(W)$.
Choose any basis of that image and any lifts $b_1,\ldots,b_h\in W$.  The resulting carriers form a
differential transcendence basis of $\K(W)$ over the field generated by the \emph{entire} tail
$W_{>N}$.  Thus each visible valuation direction contributes one complete hypertranscendental
channel, not just one algebraically independent element.

For a derivative cutoff $r$, define the ordinary intermediate field
\begin{equation}\label{eq:finite-jet-field-intro}
 L_{W,N,r}:=\K(W_{>N})
  \bigl(D^ka_{b_j}:1\le j\le h,\ 0\le k\le r\bigr).
\end{equation}
The rank is not defined as the size of a selected rectangle.  It is the intrinsic field invariant
\begin{equation}\label{eq:H-intrinsic-intro}
 H_W(N,r):=\trdeg_{\K(W_{>N})}L_{W,N,r}.
\end{equation}
The basis-and-lift invariance proved below makes this definition unambiguous.

\begin{theorem}[Valuation-quotient hypertranscendence law]\label{thm:master-intro}
For every complex subspace $W\le t\C[[t]]$ and every $N,r\ge0$, let
$h_W(N)=\dim_\C\pi_N(W)$.  Then:
\begin{enumerate}
\item the radial realization is source-faithful:
      \[
       a_u\in\K(W)\Longleftrightarrow u\in W,
       \qquad
       \K(U)\subseteq\K(W)\Longleftrightarrow U\subseteq W;
      \]
\item if $b_1,\ldots,b_{h_W(N)}$ are arbitrary lifts of an arbitrary basis of $\pi_N(W)$, then
      \[
       \K(W)=\K(W_{>N})\langle a_{b_1},\ldots,a_{b_{h_W(N)}}\rangle_D;
      \]
\item the carriers $a_{b_1},\ldots,a_{b_{h_W(N)}}$ are jointly differentially independent over
      $\K(W_{>N})$;
\item the embedded field in \emph{(2)} and the finite-jet field~\emph{\eqref{eq:finite-jet-field-intro}}
      are independent of the basis and lifts; and
\item the intrinsic finite-jet rank is
      \begin{equation}\label{eq:master-rank}
       H_W(N,r)=(r+1)h_W(N).
      \end{equation}
\end{enumerate}
Equivalently, $W/W_{>N}$ is realized as a block of $h_W(N)$ relatively
hypertranscendental channels over the complete tail.
\end{theorem}

The proof rests on a prime-profile freeness theorem for arbitrary families.  Its delicate point is
not the scalar Vandermonde matrix but the determinant in the Dirichlet-convolution ring.  We give
the missing coefficient argument explicitly: selected prime derivations make each Jacobian row
live on a distinct prime ray, and unique factorization isolates a coefficient of the convolution
determinant equal to the nonsingular scalar determinant.

There is also a natural multiplicative operation.  If $A$ is supported on positive prime powers,
write
\[
 R_p(A;t)=\sum_{n\ge1}A(p^n)t^n.
\]
Define raywise fusion by
\[
 (A\boxtimes B)(p^n)=[t^n]R_p(A;t)R_p(B;t),
\]
and set it to zero off positive prime powers.  If $\Pi_{\mathrm{pp}}$ denotes projection onto
positive-prime-power coordinates, then
\begin{equation}\label{eq:fusion-natural-intro}
 A\boxtimes B=\Pi_{\mathrm{pp}}(A*B).
\end{equation}
Thus fusion is the prime-power component of the native Dirichlet product, not an unrelated second
multiplication.  For radial sources, $a_u\boxtimes a_v=a_{uv}$.

The relation to prior work is important.  Jacobian methods for algebraic independence in
arithmetic-function fields are classical \cite{ShapiroSparer1986}, and Pong developed powerful
differential-algebraic applications of this framework \cite{Pong2016}.  Order semigroups of
subalgebras of $\C[[t]]$ are likewise established objects
\cite{AssiGarciaSanchezMicale2017,Hamilton2020}.  In particular, the semigroup of orders is already
a standard discrete invariant in Hamilton's classification of complete subalgebras.  We therefore
make no novelty claim for the order census, support realization, or conductor theorem.  The
contribution is the source-faithful radial differential-field representation and the exact
tail-relative transcendence theorem~\ref{thm:master-intro}, including the intrinsic finite-jet
field and its rank.

The paper is entirely algebraic.  No convergence of Dirichlet series is used, and no assertion
about the zeros of the Riemann zeta function or the Riemann hypothesis is made.

\section{Arithmetic differential fields}\label{sec:arithmetic}

Let $\A$ be the ring of arithmetic functions $a:\N_{>0}\to\C$ with pointwise addition and
Dirichlet convolution
\begin{equation}\label{eq:dirichlet-convolution}
 (a*b)(n)=\sum_{d\mid n}a(d)b(n/d).
\end{equation}
It is an integral domain \cite{CashwellEverett1959}; write $\F=\Frac(\A)$.  Define
\begin{equation}\label{eq:log-derivation}
 (Da)(n)=\mathrm i\log(n)a(n).
\end{equation}
The identity $\log(dm)=\log d+\log m$ gives the Leibniz rule for convolution, so $D$ extends
uniquely to $\F$ by the quotient rule.

For a prime $q$, define a prime-coordinate derivation on $\A$ by
\begin{equation}\label{eq:prime-derivation}
 (\partial_q a)(n)=v_q(nq)a(nq).
\end{equation}
These derivations also extend to $\F$.  Their commutator with $D$ is exact.

\begin{proposition}[Prime--Euler commutator]\label{prop:commutator}
For every prime $q$ and every $x\in\F$,
\begin{align}
 [\partial_q,D]&=\mathrm i\log(q)\,\partial_q,
                                                        \label{eq:commutator}\
 \partial_qx=0&\Longrightarrow \partial_q(Dx)=0.       \label{eq:no-activation}
\end{align}
\end{proposition}

\begin{proof}
For $a\in\A$, evaluation at $n$ gives
\begin{align*}
 (\partial_qDa-D\partial_qa)(n)
 &=v_q(nq)\bigl(\mathrm i\log(nq)-\mathrm i\log n\bigr)a(nq)\\
 &=\mathrm i\log(q)(\partial_qa)(n).
\end{align*}
The commutator of two derivations is a derivation, and the quotient-rule extension is determined by
its values on $\A$, so the identity extends to $\F$.  Equation~\eqref{eq:no-activation} follows
immediately.
\end{proof}

Define the active-prime set and the finite-prime coefficient field by
\begin{equation}\label{eq:Efin}
 \operatorname{Act}(x)=\{q\in\PP:\partial_qx\ne0\},
 \qquad
 \Efin=\{x\in\F:\operatorname{Act}(x)\text{ is finite}\}.
\end{equation}
The union bound for sums and products and the quotient rule show that $\Efin$ is a subfield of
$\F$.  Proposition~\ref{prop:commutator} shows that it is stable under $D$.

For $u\in\Pcal$, let $a_u$ be the source~\eqref{eq:radial-source}, viewed also in $\F$.  The map
$u\mapsto a_u$ is injective and complex-linear.  For a subspace $W\le\Pcal$, define
\begin{equation}\label{eq:KW}
 \K(W)=\Efin\langle a_u:u\in W\rangle_D.
\end{equation}
Equivalently, one may adjoin the carriers belonging to any Hamel basis of $W$.

\section{The separated-prime Jacobian}\label{sec:jacobian}

This section supplies the full determinant argument used later.  We first isolate the coefficient
law that prevents cancellation in the Dirichlet-convolution determinant.

\begin{lemma}[Separated-prime determinant coefficient]\label{lem:separated-coefficient}
Let $I$ be finite, let $(q_i)_{i\in I}$ be pairwise distinct primes, and let
$M=(M_{ij})_{i,j\in I}$ be a square matrix over $\A$.  Assume row $i$ is supported on the
nonnegative powers of $q_i$: whenever $M_{ij}(n)\ne0$, one has $n=q_i^e$ for some $e\ge0$.
Then, for every exponent vector $(e_i)_{i\in I}$,
\begin{equation}\label{eq:separated-det-coeff}
 (\det M)\!\left(\prod_{i\in I}q_i^{e_i}\right)
 =\det\bigl(M_{ij}(q_i^{e_i})\bigr)_{i,j\in I}.
\end{equation}
In particular, a nonzero scalar determinant on the right implies $\det M\ne0$ in $\A$.
\end{lemma}

\begin{proof}
Expand the determinant in the convolution ring:
\[
 \det M=\sum_{\sigma\in S_I}\operatorname{sgn}(\sigma)
          \mathop{*}_{j\in I}M_{\sigma(j),j}.
\]
Fix $\sigma$.  In a nonzero contribution to the coefficient at
$Q=\prod_iq_i^{e_i}$, the factor selected from row $\sigma(j)$ must be a power
$q_{\sigma(j)}^{f_{\sigma(j)}}$.  Unique factorization gives
\[
 \prod_iq_i^{f_i}=\prod_iq_i^{e_i}
 \quad\Longrightarrow\quad f_i=e_i\quad(i\in I).
\]
Thus the indicated coefficient of the convolution product is exactly
\[
 \prod_{j\in I}M_{\sigma(j),j}(q_{\sigma(j)}^{e_{\sigma(j)}}).
\]
Summing with permutation signs is the scalar determinant in
\eqref{eq:separated-det-coeff}.  If that scalar is nonzero, one coefficient of $\det M$ is
nonzero, so the determinant itself is nonzero.
\end{proof}

We also use two finite-dimensional facts.  First, if $u_1,\ldots,u_d$ are linearly independent
sequences, then there are indices $e_1,\ldots,e_d$ for which the evaluation matrix
\begin{equation}\label{eq:profile-minor}
 P=(u_j(e_a))_{1\le a,j\le d}
\end{equation}
is nonsingular.  This follows by choosing a basis from the span of the evaluation functionals;
repeated indices would give repeated rows, so the $e_a$ are automatically distinct.  Second, the
Vandermonde matrix $(\lambda_\ell^k)_{0\le\ell,k\le m}$ is nonsingular when the
$\lambda_\ell$ are distinct.

\begin{theorem}[Prime-profile freeness]\label{thm:profile-freeness}
Let $(u_i)_{i\in I}$ be an arbitrary family in $t\C[[t]]$.  If it is linearly independent over
$\C$, then the associated carriers are differentially independent over $\Efin$:
\begin{equation}\label{eq:profile-freeness}
 \bigl(D^ka_{u_i}\bigr)_{(i,k)\in I\times\N}
 \quad\text{is algebraically independent over }\Efin.
\end{equation}
\end{theorem}

\begin{proof}
An algebraic relation involves finitely many labels and finitely many derivatives.  It therefore
suffices to prove independence of
\[
 x_{j,k}=D^ka_{u_j},
 \qquad 1\le j\le d,\quad 0\le k\le m,
\]
for arbitrary finite $d,m$.  Choose distinct valuation indices $e_1,\ldots,e_d$ such that the
profile matrix $P$ in~\eqref{eq:profile-minor} is nonsingular.

Suppose a proposed polynomial relation has coefficients $c_1,\ldots,c_s\in\Efin$.  The union of
their active-prime sets is finite.  For each valuation row $a$, choose $m+1$ primes
\[
 q_{a,0},\ldots,q_{a,m}
\]
outside that union, with all $d(m+1)$ primes distinct.  Use the derivations
$\partial_{q_{a,\ell}}$ as the Jacobian rows.  Directly from
\eqref{eq:log-derivation}--\eqref{eq:prime-derivation},
\begin{equation}\label{eq:jacobian-scalar-entry}
 \bigl(\partial_{q_{a,\ell}}x_{j,k}\bigr)(q_{a,\ell}^{e_a})
 =(e_a+1)u_j(e_a)
   \bigl(\mathrm i(e_a+1)\log q_{a,\ell}\bigr)^k.
\end{equation}
Moreover, the whole arithmetic function
$\partial_{q_{a,\ell}}x_{j,k}$ is supported on the nonnegative powers of the single prime
$q_{a,\ell}$.

The scalar matrix in~\eqref{eq:jacobian-scalar-entry}, with rows $(a,\ell)$ and columns $(j,k)$,
is nonsingular.  Indeed, suppose coefficients $C_{j,k}$ lie in its kernel.  For fixed $a$, the
$m+1$ row equations say that the polynomial
\[
 \sum_{k=0}^m\left(\sum_{j=1}^dP_{a,j}C_{j,k}\right)T^k
\]
vanishes at the $m+1$ distinct numbers
$T=\mathrm i(e_a+1)\log q_{a,\ell}$.  The corresponding Vandermonde matrix is nonsingular, hence
$\sum_jP_{a,j}C_{j,k}=0$ for every $a,k$.  Since $P$ is nonsingular, every $C_{j,k}=0$.

Let $J$ be the arithmetic-function Jacobian
\[
 J_{(a,\ell),(j,k)}=\partial_{q_{a,\ell}}x_{j,k}.
\]
Apply Lemma~\ref{lem:separated-coefficient} with exponent $e_a$ on row $(a,\ell)$.  The coefficient
of $\det J$ at
\[
 \prod_{a=1}^d\prod_{\ell=0}^m q_{a,\ell}^{e_a}
\]
is precisely the nonsingular scalar determinant just obtained.  Therefore $\det J\ne0$ in
$\A$, and hence in $\F$.

The Jacobian criterion now gives algebraic independence of the $x_{j,k}$ over the common kernel
of the selected derivations.  Every coefficient $c_\nu$ of the proposed relation belongs to that
kernel because all selected primes were chosen outside its active-prime set.  The proposed
relation is therefore zero.  Since every algebraic relation has finite support, the arbitrary
family~\eqref{eq:profile-freeness} is independent.
\end{proof}

\begin{remark}
The scalar profile minor and the logarithmic Vandermonde matrices alone do not prove the theorem.
The essential ring-valued step is Lemma~\ref{lem:separated-coefficient}, which exhibits a
specific nonzero coefficient of the convolution determinant.
\end{remark}

\section{Source recovery and the field lattice}\label{sec:source-recovery}

\begin{lemma}[Relative block lemma]\label{lem:relative-block}
Let $E\subseteq L$ be differential fields.  If the combined family
$(x_i)_{i\in I}\sqcup(y_j)_{j\in J}$ is differentially independent over $E$, then
$(y_j)_{j\in J}$ is differentially independent over
$E\langle x_i:i\in I\rangle_D$.
\end{lemma}

\begin{proof}
The full jet families $X=(D^kx_i)_{i,k}$ and $Y=(D^ky_j)_{j,k}$ are algebraically independent over
$E$.  Hence $Y$ is algebraically independent over the rational field $E(X)$.  That field is
exactly $E\langle x_i:i\in I\rangle_D$.
\end{proof}

\begin{theorem}[Exact radial source recovery]\label{thm:source-recovery}
For every complex subspace $W\le\Pcal$ and every $u\in\Pcal$,
\begin{equation}\label{eq:source-recovery}
 a_u\in\K(W)\quad\Longleftrightarrow\quad u\in W.
\end{equation}
Consequently, for arbitrary subspaces $U,W\le\Pcal$,
\begin{align}
 \K(U)\subseteq\K(W)&\quad\Longleftrightarrow\quad U\subseteq W,
                                                        \label{eq:field-order}\
 \K(U)=\K(W)&\quad\Longleftrightarrow\quad U=W.        \label{eq:field-equality}
\end{align}
\end{theorem}

\begin{proof}
If $u\in W$, membership follows from the definition and linearity.  Suppose $u\notin W$, and let
$(b_i)_{i\in I}$ be a Hamel basis of $W$.  Then $(b_i)_{i\in I}\sqcup(u)$ is linearly independent.
By Theorem~\ref{thm:profile-freeness}, the corresponding carrier family is differentially
independent over $\Efin$.  The relative block lemma says that $a_u$ is differentially
transcendental over
\[
 \Efin\langle a_{b_i}:i\in I\rangle_D=\K(W),
\]
so $a_u$ cannot belong to that field.  This proves~\eqref{eq:source-recovery}.

If $\K(U)\subseteq\K(W)$ and $u\in U$, then $a_u\in\K(W)$, so
\eqref{eq:source-recovery} gives $u\in W$.  The converse follows from differential adjunction.
This proves~\eqref{eq:field-order}; applying it in both directions gives
\eqref{eq:field-equality}.
\end{proof}

Thus the full profile subspace, not merely its dimension or order set, is recovered by the embedded
differential field.

\section{The complete-tail extension}\label{sec:tail}

Fix a subspace $W\le\Pcal$ and $N\ge0$.  Define
\begin{align}
 \pi_N:W&\longrightarrow\C^N,
 &\pi_N(u)&=([t]u,\ldots,[t^N]u),                         \label{eq:prefix-map}\
 W_{>N}&=W\cap\ker\pi_N,
 &V_N(W)&=\pi_N(W).                                      \label{eq:tail-space}
\end{align}
The space $V_N(W)$ is finite-dimensional; put $h_W(N)=\dim_\C V_N(W)$.

Choose an arbitrary basis $(\bar b_j)_{j\in J_N}$ of $V_N(W)$ and arbitrary lifts
$b_j\in W$ satisfying $\pi_N(b_j)=\bar b_j$.  Choose also a Hamel basis $(c_i)_{i\in I_N}$ of
$W_{>N}$.

\begin{lemma}[Tail--quotient splitting]\label{lem:tail-splitting}
The combined family
\begin{equation}\label{eq:combined-basis}
 (c_i)_{i\in I_N}\sqcup(b_j)_{j\in J_N}
\end{equation}
is a Hamel basis of $W$.
\end{lemma}

\begin{proof}
It is independent because the tail span lies in $\ker\pi_N$ while $\pi_N$ is injective on the
span of the $b_j$.  To prove spanning, let $u\in W$.  Express $\pi_N(u)$ in the basis
$(\bar b_j)$.  Subtract the same finite linear combination of the lifts $b_j$ from $u$.  The
remainder lies in $W_{>N}$ and is therefore a finite linear combination of the $c_i$.
\end{proof}

\begin{theorem}[Complete-tail generation and freeness]\label{thm:tail-generation}
With the notation above,
\begin{align}
 \K(W)&=\K(W_{>N})\langle a_{b_j}:j\in J_N\rangle_D,
                                                        \label{eq:tail-generation}\
 (a_{b_j})_{j\in J_N}&\text{ is differentially independent over }\K(W_{>N}).
                                                        \label{eq:tail-relative-free}
\end{align}
Consequently the lifted quotient channels form a differential transcendence basis of
$\K(W)/\K(W_{>N})$, and
\begin{equation}\label{eq:dtrdeg}
 \dtrdeg_{\K(W_{>N})}\K(W)=\dim_\C V_N(W).
\end{equation}
\end{theorem}

\begin{proof}
By Lemma~\ref{lem:tail-splitting}, adjoining the carriers of the combined family generates
$\K(W)$.  The left block generates $\K(W_{>N})$, which proves~\eqref{eq:tail-generation}.

The combined profile family is linearly independent, so its carrier family is differentially
independent over $\Efin$ by Theorem~\ref{thm:profile-freeness}.  Apply the relative block lemma
with the tail carriers as the left block.  This proves~\eqref{eq:tail-relative-free}.  The
differential-generation and differential-independence statements together say precisely that the
quotient carriers are a differential transcendence basis.  Its cardinality is
$\dim_\C V_N(W)$, giving~\eqref{eq:dtrdeg}.
\end{proof}

The base in Theorem~\ref{thm:tail-generation} is the field generated by every invisible profile,
not by a finite tail approximation.  The theorem therefore measures the exact extension across
the valuation boundary.

\section{Intrinsic finite-jet rank}\label{sec:jet-rank}

For $r\ge0$, form the ordinary intermediate field
\begin{equation}\label{eq:finite-jet-field}
 L_{W,N,r}=\K(W_{>N})
 \bigl(D^ka_{b_j}:j\in J_N,\ 0\le k\le r\bigr).
\end{equation}

\begin{theorem}[Intrinsic relative jet-rank law]\label{thm:jet-rank}
The field $L_{W,N,r}$ is independent of the chosen basis of $V_N(W)$ and of all choices of lifts.
Its ordinary transcendence degree is
\begin{equation}\label{eq:jet-rank}
 \trdeg_{\K(W_{>N})}L_{W,N,r}
   =(r+1)\dim_\C V_N(W).
\end{equation}
\end{theorem}

\begin{proof}
By~\eqref{eq:tail-relative-free}, the family
\[
 \bigl(D^ka_{b_j}\bigr)_{(j,k)\in J_N\times\{0,\ldots,r\}}
\]
is algebraically independent over $\K(W_{>N})$.  It generates the ordinary field
$L_{W,N,r}$, hence it is a transcendence basis of that field.  Its size is
$(r+1)|J_N|=(r+1)\dim_\C V_N(W)$, proving~\eqref{eq:jet-rank}.

It remains to prove choice independence.  Let $b'_1,\ldots,b'_h$ be another lifted quotient
basis.  There is an invertible matrix $C=(C_{\ell j})\in\operatorname{GL}_h(\C)$ such that
\[
 \pi_N(b'_\ell)=\sum_jC_{\ell j}\pi_N(b_j).
\]
Thus
\[
 d_\ell:=b'_\ell-\sum_jC_{\ell j}b_j\in W_{>N}.
\]
Linearity of the source map and of $D$ gives, for every $0\le k\le r$,
\begin{equation}\label{eq:jet-change-basis}
 D^ka_{b'_\ell}=\sum_jC_{\ell j}D^ka_{b_j}+D^ka_{d_\ell}.
\end{equation}
The last term belongs to $\K(W_{>N})$.  Hence the primed finite-jet field is contained in the
unprimed field.  Applying the same argument to $C^{-1}$ gives the reverse inclusion.  This also
covers a change of lifts with the quotient basis fixed, for which $C$ is the identity.
\end{proof}

We therefore define the \emph{relative jet-rank profile} intrinsically by
\begin{equation}\label{eq:H-definition}
 H_W(N,r):=\trdeg_{\K(W_{>N})}L_{W,N,r}.
\end{equation}
Theorem~\ref{thm:jet-rank} gives $H_W(N,r)=(r+1)h_W(N)$ and completes the proof of
Theorem~\ref{thm:master-intro}.  We avoid the phrase ``Kolchin rank surface'': the invariant in
\eqref{eq:H-definition} is an ordinary transcendence degree of a bounded jet field, while
\eqref{eq:dtrdeg} is the corresponding differential transcendence degree of the full extension.

\section{The order filtration as a classical shadow}\label{sec:order-shadow}

For a subspace $W\le\Pcal$, define
\[
 \Omega(W)=\{\ord_tu:0\ne u\in W\}.
\]
The following filtered-vector-space identity is elementary but useful for reading
Theorem~\ref{thm:jet-rank}.

\begin{proposition}[Order census]\label{prop:order-census}
For every $N\ge0$,
\begin{align}
 h_W(N)&=\#\bigl(\Omega(W)\cap[1,N]\bigr),                \label{eq:order-census}\
 h_W(N+1)-h_W(N)&=\mathbf 1_{\Omega(W)}(N+1).             \label{eq:order-jump}
\end{align}
Consequently
\begin{equation}\label{eq:jet-order-census}
 H_W(N,r)=(r+1)\#\bigl(\Omega(W)\cap[1,N]\bigr).
\end{equation}
\end{proposition}

\begin{proof}
Deleting the last coordinate gives a surjection $V_{N+1}(W)\twoheadrightarrow V_N(W)$.  Its
kernel embeds in the one-dimensional last-coordinate space.  The kernel is nonzero precisely
when a profile has its first nonzero coefficient in degree $N+1$.  Rank--nullity gives
\eqref{eq:order-jump}; summing from $0$ gives~\eqref{eq:order-census}.  Equation
\eqref{eq:jet-order-census} follows from Theorem~\ref{thm:jet-rank}.
\end{proof}

If $B\le\Pcal$ is a nonunital subalgebra, then
\[
 \Gamma(B)=\{0\}\cup\Omega(B)
\]
is an additive submonoid of $\N$, since $\ord_t(uv)=\ord_tu+\ord_tv$.  The use of such order
semigroups in classifying subalgebras of one-variable power-series rings is classical; see
Assi--Garc\'{i}a-S\'{a}nchez--Micale \cite{AssiGarciaSanchezMicale2017} and, for complete complex
subalgebras and their moduli, Hamilton \cite{Hamilton2020}.  Likewise, if $d=\gcd\Gamma(B)>0$,
the eventual conductor law for $\Gamma(B)/d$ is standard numerical-semigroup theory
\cite{RosalesGarciaSanchez2009}.  Combined with~\eqref{eq:jet-order-census}, it gives an eventual
periodic jump formula for $H_B$.  These semigroup statements are consequences and contextual
shadows; they are not the main theorem of this paper.

\section{Fusion as projected Dirichlet convolution}\label{sec:fusion}

Define the positive-prime-power projection $\Pi_{\mathrm{pp}}:\A\to\A$ by retaining the values at
$p^n$, $n\ge1$, and setting all other coordinates to zero.  For arbitrary sources $A,B$, define
raywise Cauchy fusion by
\begin{equation}\label{eq:fusion-definition}
 (A\boxtimes B)(p^n)=\sum_{i=1}^{n-1}A(p^i)B(p^{n-i}),
 \qquad n\ge1,
\end{equation}
and set it to zero away from positive prime powers.  Empty sums are zero.

\begin{proposition}[Dirichlet-projection law]\label{prop:fusion-natural}
If $A$ and $B$ are supported on positive prime powers, then
\begin{equation}\label{eq:fusion-natural}
 A\boxtimes B=\Pi_{\mathrm{pp}}(A*B).
\end{equation}
On this support class, $\boxtimes$ is complex-bilinear, commutative, and associative.
\end{proposition}

\begin{proof}
Support on positive prime powers implies $A(1)=B(1)=0$.  Every divisor pair of $p^n$ is uniquely
of the form $(p^i,p^{n-i})$, $0\le i\le n$.  Therefore
\[
 (A*B)(p^n)=\sum_{i=0}^nA(p^i)B(p^{n-i})
             =\sum_{i=1}^{n-1}A(p^i)B(p^{n-i}),
\]
which proves~\eqref{eq:fusion-natural}.  Bilinearity and commutativity follow either from this
formula or from multiplication of ray series.  The ray series of $A\boxtimes B$ is
$R_p(A;t)R_p(B;t)$; associativity follows from associativity of formal power-series multiplication.
\end{proof}

For the radial source~\eqref{eq:radial-source}, $R_p(a_u;t)=u$ for every prime $p$.  Hence:

\begin{corollary}[Radial fusion law]\label{cor:radial-fusion}
For all $u,v\in\Pcal$,
\begin{equation}\label{eq:radial-fusion}
 a_u\boxtimes a_v=a_{uv}.
\end{equation}
\end{corollary}

The operation is external: it is defined on arithmetic sources and is not asserted to be a
second multiplication on all of $\F$.  We call $\K(W)$ \emph{radially fusion closed} when
$a_u,a_v\in\K(W)$ implies $a_u\boxtimes a_v\in\K(W)$ for radial profiles $u,v$.

\begin{theorem}[Fusion-closure classifier]\label{thm:fusion-classifier}
The field $\K(W)$ is radially fusion closed if and only if $W$ is closed under multiplication.
Consequently
\begin{equation}\label{eq:fusion-order-embedding}
 B\longmapsto\K(B)
\end{equation}
is an order embedding from nonunital complex subalgebras of $t\C[[t]]$ into the intermediate
differential fields of $\F/\Efin$.
\end{theorem}

\begin{proof}
Suppose $\K(W)$ is radially fusion closed and $u,v\in W$.  Source recovery gives
$a_u,a_v\in\K(W)$, so~\eqref{eq:radial-fusion} gives $a_{uv}\in\K(W)$.  A second application of
source recovery gives $uv\in W$.  The converse is the same argument in reverse.  The
order-embedding assertion is~\eqref{eq:field-order}--\eqref{eq:field-equality} restricted to
multiplicatively closed subspaces.
\end{proof}

\section{Literature boundary and contribution}\label{sec:literature}

Three parent theories meet in the result.

\paragraph{Arithmetic-function Jacobians.}
Shapiro and Sparer developed algebraic-independence criteria for Dirichlet series from derivations
of arithmetic-function fields \cite{ShapiroSparer1986}.  Quantitative variants were studied by
Komatsu, Laohakosol, and Ruengsinsub \cite{KomatsuLaohakosolRuengsinsub2011}.  Pong explains and
extends the differential-algebraic setting, including the logarithmic derivation and strong
single-source applications \cite{Pong2016}.  The Jacobian criterion and the idea of prime
derivations are therefore parent methods.  The separated-ray coefficient extraction in
Lemma~\ref{lem:separated-coefficient} is included to make the present simultaneous-family proof
complete, not as an isolated priority claim.

\paragraph{Differential algebra.}
The relative block argument, differential transcendence bases, and ordinary transcendence degree
of a rational field generated by an algebraically independent family are standard differential
and field algebra; see Kolchin \cite{Kolchin1973}.  The contribution is their exact realization
by prime-radial source spaces, with $W/W_{>N}$ controlling the extension over the complete tail.

\paragraph{Power-series order semigroups.}
Order semigroups and standard bases for subalgebras of $K[[t]]$ are established
\cite{AssiGarciaSanchezMicale2017}.  Hamilton treats the semigroup of orders as the discrete
invariant underlying moduli of complete complex subalgebras \cite{Hamilton2020}.  Thus the
additive order monoid, its support-algebra realization, cumulative order counts, and conductor
periodicity are not claimed as new.

Against these parent results, the specific theorem surface established here is:
\begin{enumerate}
\item an arbitrary-cardinality, source-faithful embedding of prime-radial profile subspaces into
      one fixed arithmetic differential-field lattice;
\item exact membership, containment, and equality of the resulting embedded fields;
\item generation of $\K(W)$ over the complete valuation tail by a basis of the finite quotient
      $W/W_{>N}$;
\item relative differential independence of that quotient basis; and
\item basis- and lift-independent finite-jet fields whose transcendence degrees are exactly
      $(r+1)\dim_\C\pi_N(W)$.
\end{enumerate}
To the best of the author's knowledge, the cited literature does not state this simultaneous
tail-relative differential-field classification.  This is a scoped novelty statement, not a
claim that its parent tools or its semigroup shadow are new.

\section{Scope and limitations}\label{sec:scope}

The results concern formal arithmetic functions and their fraction field.  They do not by
themselves produce analytic functions on a half-plane, and no convergence, meromorphic
continuation, functional equation, positivity property, or zero-set statement is used.

The radial hypothesis is also substantive: the coefficient at $p^n$ depends only on $n$.  A
prime-dependent family may carry information invisible to the one-variable profile space and is
not classified here.  Fusion is used only on the positive-prime-power carrier locus and is not
promoted to a multiplication on arbitrary elements of the fraction field.

Finally, the theorem does not imply the Riemann hypothesis.  It supplies a structural boundary:
within the radial sector, all finite-cutoff differential-field information is exhausted by the
valuation quotient $W/W_{>N}$.  Any application to zeta would first require a separate analytic
bridge and, if it needs information beyond this quotient geometry, genuinely prime-dependent,
archimedean, reflection, positivity, or completion data.

\section{Conclusion}

Prime-radial sources turn valuation profiles into concrete elements of one arithmetic
differential field.  Prime separation makes that realization faithful: no nonlinear differential
relation identifies distinct profile subspaces.  Across a valuation cutoff, the result is sharper
than a dimension count.  The complete invisible tail is an exact differential base, the visible
quotient supplies a differential transcendence basis, and every bounded jet block has an
intrinsic ordinary transcendence degree equal to its expected product dimension.

The multiplicative picture is compatible with the native arithmetic product: raywise fusion is
Dirichlet convolution followed by positive-prime-power projection.  In the radial sector this
reduces to multiplication in $t\C[[t]]$, while classical order semigroups record only the filtered
shadow of the stronger field correspondence.  The central law is therefore the tail-relative
one: every finite valuation quotient is realized, without loss or hidden collapse, as a block of
relatively hypertranscendental channels.

\section*{Declaration of generative AI and AI-assisted technologies in the manuscript preparation process}

During the preparation of this work, the author used GPT 5.6 Sol, Fable 5, and Grok 4.5 for
research assistance, literature organization, proof-review support, and manuscript drafting.
The author reviewed and edited the resulting material, verified the mathematical claims and
references, and takes full responsibility for the content of the article.

\begin{thebibliography}{99}

\bibitem{AssiGarciaSanchezMicale2017}
A.~Assi, P.~A. Garc\'{i}a-S\'{a}nchez, and V.~Micale,
\newblock Bases of subalgebras of $K[[x]]$ and $K[x]$,
\newblock \emph{J. Symbolic Comput.} \textbf{79} (2017), 414--425,
\newblock \url{https://doi.org/10.1016/j.jsc.2016.08.002}.

\bibitem{CashwellEverett1959}
E.~D. Cashwell and C.~J. Everett,
\newblock The ring of number-theoretic functions,
\newblock \emph{Pacific J. Math.} \textbf{9} (1959), 975--985,
\newblock \url{https://doi.org/10.2140/pjm.1959.9.975}.

\bibitem{Hamilton2020}
E.~Hamilton,
\newblock Classifying complete $\C$-subalgebras of $\C[[t]]$,
\newblock \emph{Manuscripta Math.} \textbf{163} (2020), 437--462,
\newblock \url{https://doi.org/10.1007/s00229-019-01162-5}.

\bibitem{Kolchin1973}
E.~R. Kolchin,
\newblock \emph{Differential Algebra and Algebraic Groups},
\newblock Pure and Applied Mathematics, vol.~54, Academic Press, New York, 1973.

\bibitem{KomatsuLaohakosolRuengsinsub2011}
T.~Komatsu, V.~Laohakosol, and P.~Ruengsinsub,
\newblock Independence measures of arithmetic functions,
\newblock \emph{J. Number Theory} \textbf{131} (2011), 1--17,
\newblock \url{https://doi.org/10.1016/j.jnt.2010.06.011}.

\bibitem{Pong2016}
W.~Y. Pong,
\newblock Applications of differential algebra to algebraic independence of arithmetic
functions,
\newblock \emph{Acta Arith.} \textbf{172} (2016), 149--173,
\newblock \url{https://doi.org/10.4064/aa8112-12-2015}.

\bibitem{RosalesGarciaSanchez2009}
J.~C. Rosales and P.~A. Garc\'{i}a-S\'{a}nchez,
\newblock \emph{Numerical Semigroups},
\newblock Developments in Mathematics, vol.~20, Springer, New York, 2009,
\newblock \url{https://doi.org/10.1007/978-1-4419-0160-6}.

\bibitem{ShapiroSparer1986}
H.~N. Shapiro and G.~H. Sparer,
\newblock On algebraic independence of Dirichlet series,
\newblock \emph{Comm. Pure Appl. Math.} \textbf{39} (1986), 695--745,
\newblock \url{https://doi.org/10.1002/cpa.3160390602}.

\end{thebibliography}

\end{document}
